% RESEARCH_DRAFT_SHA256: 3b270429ebb6bd0ff43e84514ff8e6e0ec6173b7fb7ba57716325ab20da2dfba
% Layout and source-reference wrappers expanded for the web renderer.
\documentclass[11pt]{article}
\usepackage[margin=25mm]{geometry}
\usepackage{amsthm}
\usepackage{fontspec,unicode-math}
\setmainfont{Latin Modern Roman}
\setmathfont{Latin Modern Math}
\usepackage{amsmath,amssymb,mathtools,xurl,hyperref}
\hypersetup{colorlinks=true,urlcolor=blue,linkcolor=blue}
\setcounter{secnumdepth}{0}
\setlength{\parindent}{0pt}
\setlength{\parskip}{5pt}
\setlength{\emergencystretch}{3em}
\newcommand{\sref}[2]{\hyperlink{source:#1:#2}{[S#2]}}
\newcommand{\cataloguescope}[1]{\paragraph{Source record.}#1}
\newcommand{\floor}[1]{\left\lfloor#1\right\rfloor}
\newtheorem{proposition}{Proposition}
\newtheorem{lemma}[proposition]{Lemma}
\begin{document}
\paragraph{Research manuscript.} OpenAI \texttt{gpt-6-astra}, reasoning \texttt{ultra}; 11 October 2026.
\paragraph{Preserved source body.} The complete catalogue mathematical text and sources follow with layout and source links expanded, including source qualifications and its historical preparation notes. New research begins at the section headed Research attempt.

% UnsolvedMath ID 20002198; source catalogue code AIM-INFRASTRUCTURE-0098
\section{A provenance-bounded design for the history of women in combinatorics}\label{20002198}
{\small UnsolvedMath ID 20002198 · Computer Science}
\subsection{Definitions and mathematical statement}

\paragraph{Shared notation.}Standard number systems are $\mathbb N,\mathbb Z,\mathbb Q,\mathbb R,\mathbb C$. Specialized notation and index ranges follow the cited source and the definitions below.


A historical study documents people, institutions, events, and their context using evidence. A demographic dataset associates recorded characteristics with a specified population and time period. Counts and rates require a denominator and consistent inclusion rules; accurate gender data requires a defined collection procedure rather than inference from names.
\paragraph{Record scope and qualifications.}
The catalogue title is retained as metadata. Named certificates, contracts, bounds or frameworks in that title are not asserted as results of the original source extract.
The source combines historical, sociological, and statistical questions without a fixed population or study protocol. No gender classifications, departmental ranking, or causal conclusion are assumed.

Catalogue screening limitation: Community, history, equity, or organizational policy discussion, rather than a mathematical statement.

The source snapshot labels this record partially solved; this is catalogue metadata rather than independent verification of a mathematical result or the current completion of a project.
\paragraph{Statement recorded in the supplied source.}
The History of Women in Combinatorics

What is the history of women in combinatorics? Answering this question might involve researchers from many different areas, including historians, sociologists, mathematicians, and data scientists. As other possible guiding questions: What departments are producing the most women PhDs in math? What are they doing right? What is going wrong at other departments? How can we ensure our gender data is accurate for this study? Will this study be US-centered, and if so, how can we perform similar studies outside the US? \href{http://aimpl.org/gemscombin/7/}{[S2]}
\paragraph{Source status.}
The supplied catalogue labels this source record partially solved; that metadata was checked against the supplied snapshot on 2026-10-11. The label is preserved for provenance and is not an independent verification of present open status or mathematical correctness.
\subsection{Short English statement}
Investigate the history and institutional participation of women in combinatorics while addressing data quality and geographic scope.
\subsection{Sources}
\begin{itemize}
\item[S1] Ulam AI and the UnsolvedMath Contributors, UnsolvedMath: A Curated Collection of Open Mathematics Problems, record 20002198 (AIM-INFRASTRUCTURE-0098). Catalogue curation and metadata: CC BY 4.0. The source extract is retained; definitions, English summaries and qualifications are editorial additions.
\url{https://huggingface.co/datasets/ulamai/UnsolvedMath}.
\item[S2] AIM Workshop Problem Lists, source linked by catalogue record AIM-INFRASTRUCTURE-0098.
\url{http://aimpl.org/gemscombin/7/}.
\end{itemize}
\paragraph{Preparation and review.}
The original source and UnsolvedMath contributors retain attribution for the source extract. Codex (GPT-6; exact serving revision not exposed) assisted with definitions, English summaries, scope notes and formatting. No new human mathematical authorship or independently verified solution is asserted. Permanent source ID 20002198 and code AIM-INFRASTRUCTURE-0098 checked against the complete supplied registry and current repository. Distinct upstream records, including repeated questions and broad or damaged extracts, are retained for complete coverage of the source CS category. A repeated record is not a new independent conjecture. Complete mathematical-equivalence and current-literature review has not been performed. Preparation and final review follow contribution prompts C and D, with source gaps disclosed. Structural and rendering checks do not establish a new theorem or open status.
\label{end:20002198}

\section{Research attempt}
\subsection{Literature, source context and scope}
The complete catalogue body above retains the original agenda, source status, and historical preparation. The AIMPL problem page could not be retrieved. The supplied complete source transcription was read, and the accessible primary AIM workshop announcement was checked for context. The announcement is not a substitute for verification of the inaccessible problem page. The catalogue's earlier proposed frameworks and research summaries are context, not established solutions credited to this attempt. The following is an explicitly limited investigation of one route; it neither substitutes a formalized toy model for the entire source agenda nor claims empirical or deployed-software results. Elementary auxiliary facts are not new original-target progress.

\subsection{Testing department comparisons with incomplete demographic records}
The historical question is not answered by a counting lemma. To test one proposed quantitative route, fix a particular voluntarily reported category and a cohort of $N$ people. Let $W$ be known members of that category, $K$ known nonmembers, and $M$ unrecorded cases, with $W+K+M=N$. The indicator is membership in one specified category; it does not impose an exhaustive gender binary or authorize inference of anyone's identity.
\begin{proposition}
Without assumptions on the missing records, the category proportion lies in the sharp interval $[W/N,(W+M)/N]$.
\end{proposition}
\begin{proof}
If $u$ of the $M$ missing records belong to the category, the proportion is $(W+u)/N$ with $0\le u\le M$. This gives the interval, and every integer $u$ is compatible with the observed counts by assigning that many unspecified indicator values to one and the remainder to zero. These are mathematical completions, not assignments of identities to real people.
\end{proof}
For cohort $A$ take $(N,W,M)=(100,20,40)$, giving $[1/5,3/5]$. For cohort $B$ take $(100,40,0)$, giving exactly $2/5$. Depending on the missing records, $A$ can have a lower or higher proportion than $B$. The observed counts therefore do not identify the ranking, even before comparing cohort sizes, fields, dates, definitions or causal explanations for outcomes.
\subsection{Verification and remaining gaps}
All 41 possible missing-member counts for cohort $A$ were enumerated and the endpoints and rank reversal verified. No historical archive, identity data or department outcomes were collected. A higher proportion is not by itself proof of what a department does right. The history, international comparison, data accuracy and institutional explanation questions remain unresolved.

\section{References and source attribution}
Primary source: \url{http://aimpl.org/gemscombin/7/}. Attempted primary AIMPL page retrieval; failed with timeout or gateway error. The supplied full source transcription was read; the inaccessible live page was not independently inspected.

\url{https://aimath.org/pastworkshops/gemscombin.html}. Read the primary March 2023 workshop announcement and its community inclusion agenda.

All retrievals were checked on 11 October 2026. Original contributor and catalogue attribution remains in the preserved Sources section. The elementary arguments below are self-contained; their exposition is not a claim of first discovery.


\paragraph{Integration review.} Codex (GPT-6; exact serving revision not exposed) reviewed the written arguments and expanded layout and source-reference wrappers. Repository metadata and the production web renderer were checked using Python 3.13 and JavaScript with MathJax 3.2.2. This model-assisted review is not independent human proof verification; no PDF compilation or formal proof checking is asserted.

\section{Completion Estimate}
\noindent\textbf{COMPLETION: 0\%}
No substantiated new progress toward the entire original target is claimed. The worked results and failed route are delimited above; a known result, illustrative model, or auxiliary obstruction is not counted as original-target completion.

\section{Final status and attribution}
\textbf{Attempt status: unresolved.} This describes this research attempt, not a new ruling on the source catalogue's status or on external literature claims.
The mathematics and this self-audit were generated by OpenAI \texttt{gpt-6-astra}, reasoning \texttt{ultra}, on 11 October 2026. Codex (GPT-6; exact serving revision not exposed) orchestrated the batch and contributed formatting/repository review. No human mathematical authorship, independent human proof checking, or formal verification is claimed. The historical preparation paragraphs in the reproduced catalogue body concern that body only.

Contribution prompt A was applied to the complete definition and supplied source record, followed by the adversarial checks of prompt B. Local reproduction files are \texttt{batch013/checks.py}, \texttt{computations.json}, and the per-ID audit JSON. They were not public links at preparation time. Numerical checks are finite and do not prove unrestricted results. No LaTeX engine was found locally; compilation was not performed by this research agent.

\end{document}
